\section{Условие задачи}

\subsection{Цель работы}

Написать одну законченную консольную программу --- интерактивный справочник,
работающий с динамической структурой данных заданного варианта и выполняющий
над ней набор операций, выбираемых из текстового меню. Вся логика оформляется
как библиотека чистых функций: функции не вводят и не выводят данные сами, а
принимают их параметрами, изменяют переданные структуры (в том числе по ссылке
через out-параметры) и возвращают результат. Работа состоит из двух частей:

\begin{enumerate}
  \item Единая программа варианта: динамическая структура данных, базовые
        операции семейства, алгоритм варианта и поиск с возвратом результата и
        найденного значения по ссылке.
  \item Проект CMake с Google Test, который проверяет чистые функции программы.
\end{enumerate}

\subsection{Семейство A: линейный динамический массив}

Вариант закреплён за одним из трёх семейств структур данных. Варианты 1--8
относятся к семейству A --- линейному динамическому массиву. Базовые операции
реализуются в \texttt{array\_ops.h} / \texttt{array\_ops.cpp}: создание,
освобождение, изменение размера, вставка элемента, удаление элемента, печать.

\subsection{Вариант 2}

Сортировка выбором по возрастанию; затем
\texttt{void array\_reverse(int* arr, size\_t size)} --- развернуть массив
на месте двумя указателями. Программа выводит массив после сортировки и после
разворота.

Контрольный пример из задания:
\[
  \{3,\,1,\,4,\,1,\,5\} \;\longrightarrow\; \{1,\,1,\,3,\,4,\,5\}
  \;\longrightarrow\; \{5,\,4,\,3,\,1,\,1\}.
\]

\subsection{Обязательства программы}

\begin{itemize}
  \item Никаких утечек памяти: каждая созданная через \texttt{new[]} структура
        освобождается через \texttt{delete[]}, после чего указатель
        обнуляется: \texttt{delete[] arr;\allowbreak arr = nullptr;}.
  \item Корректный ввод: при вводе чисел и меню проверяется успешность
        \texttt{std::cin} с последующим \texttt{std::cin.clear()}, программа не
        падает при некорректном вводе.
  \item Устойчивость к пустой структуре: функции поиска и алгоритмы корректно
        работают с массивом длины 0 и возвращают \texttt{false} либо безопасное
        значение.
  \item Каждая операция выводит понятный результат: сам массив, длину,
        найденные значения.
\end{itemize}

\subsection{Запрещённые конструкции}

Запрещены:

\begin{itemize}
  \item классы, конструкторы, семантика копирования и перемещения;
  \item \texttt{std::string}, \texttt{std::vector}, \texttt{std::array};
  \item умные указатели.
\end{itemize}

\noindent Работа с памятью --- только вручную (\texttt{new[]} /
\texttt{delete[]}), строки --- только C-строки.

\subsection{Часть 2: сборка и тестирование}

Требуется проект CMake и не менее 10 тестов Google Test, проверяющих создание
и удаление массива, вставку в середину, удаление и алгоритм варианта на
контрольном примере. Общие параметры работы: язык C++20, система сборки
CMake 3.16+, тестирование --- Google Test.
